\section{Conclusion and Future Work}
\label{sec:conclusion}

This paper has presented \emph{emergence-study}, an artwork in which a simulated world built from the rules of everyday life is never shown, and is seen only through software-interpretations: thoughts written in English from lived experience and realised as algorithms that read the world. I have described the base system in full, the three-part format of thought, expression and parameters, and six interpretations made with it, and I have reported how measuring and revising the base system changed what those interpretations could show.

Two directions follow. The first is to open the system to other readers. Because the interface is narrow and the format is plain, other artists can write interpretations of the same world, and a body of readings by many people would test whether one system of everyday rules can hold as many readings as I hope. The second is to keep the ground honest as it grows: to add rules only where life suggests them, and to measure what each one does to the world before reading it. Some of what such measurement uncovers, like the neighbourhoods that homes produced, may become interpretations in their own right.

\identifying{The system, the interpretations and the analysis pages are open source at \url{https://github.com/arjunmakesthings}.\note{put the exact repository link here.}}{The system, the interpretations and the analysis pages are open source; the link is withheld for review.}
